\section{Joining tables}
\subsection{Protect the unit}

\begin{frame}{A join can multiply rows quietly}
  \begin{center}
  \begin{tikzpicture}[
    table/.style={draw=palegray,rounded corners=2pt,text width=3.0cm,
      minimum height=1.55cm,align=center},
    arr/.style={-{Latex[length=2mm]},color=crimson,thick}
  ]
    \node[table] (left) {\textbf{Analysis table}\\Kenya, 2022\\one row};
    \node[table,right=1.1cm of left] (right) {\textbf{Metadata}\\Kenya\\two rows};
    \node[table,right=1.1cm of right] (joined) {\textbf{Joined table}\\Kenya, 2022\\two rows};
    \draw[arr] (left)--(right);
    \draw[arr] (right)--(joined);
  \end{tikzpicture}
  \end{center}

  \vspace{0.65em}
  Every value can look real while Kenya receives twice the weight. Inspect the
  join structure because the values alone will hide this problem.
\end{frame}

\begin{frame}[fragile]{Test the right-side key first}
\begin{lstlisting}[language=R]
stopifnot(
  !anyDuplicated(metadata["iso3c"])
)

joined <- left_join(
  health,
  metadata,
  by = "iso3c"
)
\end{lstlisting}

  The key on the right table determines whether rows on the left can multiply.
\end{frame}

\begin{frame}[fragile]{Leave a join record}
\begin{lstlisting}[language=R]
stopifnot(nrow(joined) == nrow(health))
stopifnot(!anyDuplicated(joined[c("iso3c", "year")]))

unmatched <- anti_join(
  health,
  metadata,
  by = "iso3c"
)
\end{lstlisting}

  Successful syntax still requires row-count and key checks to confirm that the
  join preserved the object described by the table contract.
\end{frame}

\begin{frame}{Diagnose the broken join}
  The row count increased after \texttt{left\_join()}.

  Which check should come first?
  \begin{enumerate}
    \item Delete duplicated rows in the output.
    \item Test the right-side key for duplicates.
    \item Change to \texttt{inner\_join()}.
    \item Ask an agent to rewrite the complete pipeline.
  \end{enumerate}

  \vfill
  The first useful check is number 2. The other options can hide the reason.
\end{frame}
